\color{blue}
\subsection{Distributed Path Sampling}\todo{inline equation if possible}
\label{sec:distrib-path}

The presence of semi-trusted third parties (STTPs) in our decentralized header construction naturally enables an additional functionality: distributing the path selection process itself. 
In classical Sphinx-based systems, the client unilaterally selects the routing path by choosing a sequence of mixnodes and deriving the corresponding shared secrets from their public keys. 
While simple, this design gives the client full control over route selection, which may enable biased routing strategies or facilitate route fingerprinting attacks.

To mitigate this issue, we propose a mechanism that distributes path selection with the assistance of STTPs. 
Instead of allowing the client to directly compute the shared secrets associated with a path, the STTPs collaboratively provide the cryptographic material required to derive them.

In the current design, this functionality is implemented as a protocol executed prior to the construction of the DSphinx packet. 
The STTPs return partial shares that allow the client to reconstruct the shared secrets associated with the sampled path. 
The client can then proceed with the DSphinx header construction procedure described in Fig.~\ref{fig:overall_schema}, as if it had derived the secrets itself.

Although effective, this mechanism remains conceptually separate from the header construction protocol. 
We therefore conclude this section by briefly discussing how distributed path sampling could be integrated more tightly into the decentralized header construction.

\subsubsection{Overview}

In standard Sphinx routing, the client derives a sequence of shared secrets $s_i$ from the public keys $K_i$ of the selected mixnodes. 
In contrast, our approach relies on the mixnodes' secret keys $k_i$, which are distributed among STTPs using Shamir Secret Sharing (SSS). 
As a result, no single entity possesses sufficient information to compute the routing secrets independently.

Consequently, the client cannot derive routing secrets on its own. 
Instead, it queries a set of STTPs, which collaboratively sample a path in a pseudorandom manner and provide shares of the material required to reconstruct the corresponding secrets.

After collecting shares from a threshold number of STTPs, the client reconstructs the shared secrets and proceeds with the DSphinx header construction procedure (Fig.~\ref{fig:overall_schema}).

\subsubsection{Setup.}

The goal of the setup phase is to distribute shares of each mixnode's secret key and routing information among the STTPs.
%
Each STTP $t$ is identified by a value $ID_t = H(IP_t)$, where $IP_t$ denotes the IP address of the STTP. 
Each mixnode $i$ samples a random identifier $r_i$ and constructs two random polynomials $f_i$ and $g_i$ of degree $n-1$ respectively satisfying the constraint $f_i(0) = k_i$, and $g_i(0) = N_i$. 
Here, $k_i$ denotes the mixnode's secret key, $N_i$ its address point, and $n$ the reconstruction threshold of the Shamir sharing scheme.

The mixnode then distributes to each STTP $t$ the share defined in Eq.~\ref{eq:sttp_share}. 
STTPs store these shares indexed by the identifier $r_i$, which later serves to pseudorandomly select mixnodes during path sampling.
Ideally, shares are delivered anonymously so that STTPs cannot associate them with specific mixnodes.
\begin{equation}
    (r_i,\quad k_{i_t}=f_i(ID_t),\quad N_{i_t}=g_i(ID_t))
    \label{eq:sttp_share}
\end{equation}

Adding a new mixnode only requires executing the setup procedure for that node. 
Adding a new STTP requires each mixnode to generate and transmit an additional share. 
If the threshold parameter $n$ changes, however, the entire setup procedure must be repeated.

\subsubsection{Path Sampling.}

To sample a routing path, the client sends a nonce $w$ to a set of $n$ STTPs. 
This nonce acts as a randomness seed. 
Each STTP first computes the elliptic curve point $\A{} = \Map(w)$, where $\Map(\cdot)$ denotes the Elligator mapping. 
This mapping produces a uniformly distributed elliptic curve point whose discrete logarithm relative to the base point $G$ is unknown.

For each hop position $i$, the STTP selects the mixnode share $(r_i, k_{i_t}, N_{i_t})$ whose identifier $r_i$ is closest to $H(w \parallel \text{timestamp} \parallel i)$.
Then it computes a share of the Diffie-Hellman value associated with the selected mixnode shares as $S_{i_t} = k_{i_t} \A{}$,
and returns to the client the tuple $(ID_t,\, N_{i_t},\, S_{i_t})$.

\subsubsection{Secret Reconstruction.} 

Upon receiving responses from a threshold number of STTPs, the client reconstructs the mixnode address points and shared secret points using Lagrange interpolation, as shown in Eq.~\ref{eq:lagrange_reconstruction}. 
Here, $\lambda(ID_t)$ denotes the Lagrange interpolation coefficient associated with STTP identifier $ID_t$. 
Since the same nonce $w$ is used by all STTPs, the reconstructed points correspond to the same sequence of mixnodes selected during the path sampling phase.
\begin{equation}
    N_i = \sum_t \lambda(ID_t) N_{i_t}, \qquad
    S_i = \sum_t \lambda(ID_t) S_{i_t},
    \label{eq:lagrange_reconstruction}
\end{equation}

To prevent linkability, the client samples a random scalar $r$ and rerandomizes the shared secrets by computing $S'_i = r S_i$. 
The corresponding cryptographic header element is updated accordingly as $\A{}' = r\A{}$.

Since DSphinx derives shared secrets sequentially, the client computes the scalar shared secrets in cascade. 
For each hop $i$, the corresponding scalar secret is obtained by combining the reconstructed point with the secrets of the previous hops and mapping the result back to the scalar field, as shown in Eq.~\ref{eq:cascade_secret}. 
\begin{equation}
    s'_i = \Map^{-1}\!\left(S'_i \prod_{j=1}^{i-1} s'_j \right).
    \label{eq:cascade_secret}
\end{equation}

Finally, the client proceeds with the standard DSphinx header construction procedure (Fig.~\ref{fig:overall_schema}) using $\A{}'$, together with the reconstructed values $N_i$ and $s'_i$ for each hop $i$ along the path.
